% --- CORRECCIÓN IMPORTANTE EN LA PRIMERA LÍNEA ---
% Agregamos "xcolor=table" para que funcione \rowcolor sin errores
\documentclass[aspectratio=169, 10pt, xcolor=table]{beamer}

% --- Configuración del Tema (estilo profesional como el ejemplo) ---
\usetheme{Madrid}
\usecolortheme{dolphin}

% --- Paquetes necesarios ---
\usepackage[utf8]{inputenc}
\usepackage[spanish]{babel}
\usepackage{amsmath, amssymb}
\usepackage{booktabs}
\usepackage{graphicx}
\usepackage{listings}
\usepackage{tikz}
\usetikzlibrary{arrows.meta, positioning, shapes.geometric}

% --- Ruta de Imágenes (crea la carpeta "imagenes" y coloca ahí tus archivos) ---
\graphicspath{{./imagenes/}}

% --- Estilo para código Python ---
\definecolor{codegreen}{rgb}{0,0.6,0}
\definecolor{codegray}{rgb}{0.5,0.5,0.5}
\definecolor{codepurple}{rgb}{0.58,0,0.82}
\definecolor{backcolour}{rgb}{0.95,0.95,0.92}
\lstset{
    language=Python,
    backgroundcolor=\color{backcolour},
    commentstyle=\color{codegreen},
    keywordstyle=\color{blue},
    numberstyle=\tiny\color{codegray},
    stringstyle=\color{codepurple},
    basicstyle=\ttfamily\scriptsize,
    breaklines=true,
    frame=single,
    showstringspaces=false
}

% --- Pie de página personalizado ---
\makeatletter

\setbeamertemplate{footline}{
  \leavevmode%
  \hbox{%
  \begin{beamercolorbox}[wd=.333333\paperwidth,ht=2.25ex,dp=1ex,center]{author in head/foot}%
    \usebeamerfont{author in head/foot}\insertshortauthor
  \end{beamercolorbox}%
  \begin{beamercolorbox}[wd=.333333\paperwidth,ht=2.25ex,dp=1ex,center]{title in head/foot}%
    \usebeamerfont{title in head/foot}Sesión 6
  \end{beamercolorbox}%
  \begin{beamercolorbox}[wd=.333333\paperwidth,ht=2.25ex,dp=1ex,right]{date in head/foot}%
    \usebeamerfont{date in head/foot}NLP \hfill \insertframenumber/\inserttotalframenumber \hspace*{0.3cm}
  \end{beamercolorbox}}%
  \vskip0pt%
}

\makeatother

% --- Datos de la portada ---
\title[SESIÓN 6 NLP]{\textbf{Sesión 6}}
\subtitle{Redes Neuronales Recurrentes (RNN, LSTM, GRU)}
\author[Prof. C. Sánchez]{Carlos César Sánchez Coronel}
\institute{\textbf{NLP}}
\date{}

% Logo opcional (descomentar si tienes la imagen)
% \titlegraphic{\includegraphics[height=1.5cm]{logo.png}}

% --- Eliminar la barra de navegación (opcional) ---
\setbeamertemplate{navigation symbols}{}

% --- Índice al inicio de cada sección ---
\AtBeginSection[]{
  \begin{frame}
    \frametitle{Contenido}
    \tableofcontents[currentsection]
  \end{frame}
}

% ============================================================================
% INICIO DEL DOCUMENTO
% ============================================================================
\begin{document}

% --- Portada ---
\begin{frame}
    \titlepage
\end{frame}

% --- Índice general ---
\begin{frame}{Contenido}
    \tableofcontents
\end{frame}

% ============================================================================
% SECCIÓN 1: INTRODUCCIÓN
% ============================================================================
\section{Introducción}

\begin{frame}{Limitaciones de las redes feed-forward para texto}
    \begin{itemize}
        \item Las redes feed-forward con average pooling pierden el orden de las palabras.
        \item No pueden distinguir: ``El producto es bueno, pero la atención es pésima'' vs ``La atención es pésima, pero el producto es bueno''.
        \item \textbf{Solución}: Redes Neuronales Recurrentes (RNN), que procesan secuencias manteniendo un estado oculto.
    \end{itemize}
    % Basado en Pregunta 1
\end{frame}

% ============================================================================
% SECCIÓN 2: RNN VANILLA
% ============================================================================
\section{RNN Vanilla}

\begin{frame}{Arquitectura de una RNN}
    \begin{block}{Ecuación de actualización}
        \[
        h_t = \tanh(W_{hh} h_{t-1} + W_{xh} x_t + b_h)
        \]
        \begin{itemize}
            \item $x_t$: entrada en el paso $t$.
            \item $h_{t-1}$: estado oculto anterior.
            \item $h_t$: nuevo estado oculto.
        \end{itemize}
    \end{block}
    % IMAGEN: Diagrama de RNN desenrollada
\end{frame}

\begin{frame}{Backpropagation Through Time (BPTT)}
    \begin{itemize}
        \item Desenrollar la RNN para cada paso de tiempo.
        \item Calcular gradientes sumando contribuciones de cada paso.
        \item Problema: gradiente puede \textbf{desvanecerse} o \textbf{explotar}.
    \end{itemize}
    % Basado en Preguntas 2 y 5
\end{frame}

\begin{frame}{Problema del gradiente desvaneciente}
    \begin{itemize}
        \item Durante BPTT, el gradiente se multiplica por $W_{hh}$ repetidamente.
        \item Si $\|W_{hh}\| < 1$, el gradiente tiende a cero exponencialmente.
        \item La red no puede aprender dependencias de largo alcance.
        \item \textbf{Soluciones}: LSTM, GRU, inicialización cuidadosa, clipping.
    \end{itemize}
    % Basado en Pregunta 2
\end{frame}

% ============================================================================
% SECCIÓN 3: LSTM (LONG SHORT-TERM MEMORY)
% ============================================================================
\section{LSTM (Long Short-Term Memory)}

\begin{frame}{Arquitectura LSTM}
    \begin{block}{Componentes}
        \begin{itemize}
            \item \textbf{Puerta de olvido} ($f_t$): qué información del estado anterior olvidar.
            \item \textbf{Puerta de entrada} ($i_t$): qué nueva información almacenar.
            \item \textbf{Celda candidata} ($\tilde{c}_t$): nueva información propuesta.
            \item \textbf{Puerta de salida} ($o_t$): qué información del estado pasar al siguiente paso.
        \end{itemize}
    \end{block}
    % Basado en Pregunta 3
\end{frame}

\begin{frame}{Ecuaciones LSTM}
    \begin{align*}
        f_t &= \sigma(W_f [h_{t-1}, x_t] + b_f) \\
        i_t &= \sigma(W_i [h_{t-1}, x_t] + b_i) \\
        \tilde{c}_t &= \tanh(W_c [h_{t-1}, x_t] + b_c) \\
        c_t &= f_t \odot c_{t-1} + i_t \odot \tilde{c}_t \\
        o_t &= \sigma(W_o [h_{t-1}, x_t] + b_o) \\
        h_t &= o_t \odot \tanh(c_t)
    \end{align*}
    % Basado en Pregunta 3
\end{frame}

\begin{frame}{Interpretación de las puertas}
    \begin{itemize}
        \item \textbf{Puerta de olvido} ($f_t$): valores cercanos a 1 = recordar, 0 = olvidar.
        \item \textbf{Puerta de entrada} ($i_t$): controla cuánto del nuevo candidato se añade.
        \item \textbf{Estado de celda} ($c_t$): memoria a largo plazo, flujo lineal (evita desvanecimiento).
        \item \textbf{Puerta de salida} ($o_t$): controla cuánto del estado de celda se expone.
    \end{itemize}
    % Basado en Pregunta 12
\end{frame}

% ============================================================================
% SECCIÓN 4: GRU (GATED RECURRENT UNIT)
% ============================================================================
\section{GRU (Gated Recurrent Unit)}

\begin{frame}{Arquitectura GRU}
    \begin{block}{Componentes}
        \begin{itemize}
            \item \textbf{Puerta de reinicio} ($r_t$): determina cuánto del estado anterior considerar.
            \item \textbf{Puerta de actualización} ($z_t$): combina olvido y entrada de LSTM.
        \end{itemize}
    \end{block}
    % Basado en Pregunta 15
\end{frame}

\begin{frame}{Ecuaciones GRU}
    \begin{align*}
        r_t &= \sigma(W_r [h_{t-1}, x_t] + b_r) \\
        z_t &= \sigma(W_z [h_{t-1}, x_t] + b_z) \\
        \tilde{h}_t &= \tanh(W_h [r_t \odot h_{t-1}, x_t] + b_h) \\
        h_t &= (1 - z_t) \odot h_{t-1} + z_t \odot \tilde{h}_t
    \end{align*}
    % Basado en Pregunta 15
\end{frame}

\begin{frame}{LSTM vs GRU}
    \begin{table}
        \centering
        \begin{tabular}{l|c|c}
            \toprule
            \textbf{Característica} & \textbf{LSTM} & \textbf{GRU} \\
            \midrule
            Número de puertas & 3 & 2 \\
            Parámetros & Más & Menos \\
            Estado de celda separado & Sí & No \\
            Velocidad de entrenamiento & Más lenta & Más rápida \\
            Propenso a sobreajuste & Más (con pocos datos) & Menos \\
            \bottomrule
        \end{tabular}
    \end{table}
    % Basado en Pregunta 4
\end{frame}

% ============================================================================
% SECCIÓN 5: RNN BIDIRECCIONALES
% ============================================================================
\section{RNN Bidireccionales}

\begin{frame}{RNN Bidireccional}
    \begin{itemize}
        \item Dos RNN: una procesa de izquierda a derecha (forward), otra de derecha a izquierda (backward).
        \item Estado final: $h_t = [\overrightarrow{h_t}; \overleftarrow{h_t}]$ (concatenación).
        \item Útil para tareas donde el contexto completo es necesario (NER, POS tagging).
    \end{itemize}
    % Basado en Pregunta 6
\end{frame}

% ============================================================================
% SECCIÓN 6: REGULARIZACIÓN EN RNN
% ============================================================================
\section{Regularización en RNN}

\begin{frame}{Técnicas de regularización}
    \begin{itemize}
        \item \textbf{Dropout}: en conexiones feed-forward y, con cuidado, en conexiones recurrentes (variational dropout).
        \item \textbf{Weight decay} (L2).
        \item \textbf{Early stopping}.
        \item \textbf{Gradient clipping} (para evitar explosión).
        \item \textbf{Recurrent dropout}: misma máscara en cada paso.
    \end{itemize}
    % Basado en Preguntas 10 y 19
\end{frame}

% ============================================================================
% SECCIÓN 7: IMPLEMENTACIÓN EN PYTORCH
% ============================================================================
\section{Implementación en PyTorch}

\begin{frame}[fragile]{RNN para clasificación de sentimiento}
    \begin{lstlisting}
import torch
import torch.nn as nn

class RNNClassifier(nn.Module):
    def __init__(self, vocab_size, embed_dim, hidden_dim, num_classes):
        super().__init__()
        self.embedding = nn.Embedding(vocab_size, embed_dim)
        self.rnn = nn.LSTM(embed_dim, hidden_dim, batch_first=True, bidirectional=True)
        self.fc = nn.Linear(hidden_dim*2, num_classes)
        self.dropout = nn.Dropout(0.5)

    def forward(self, x):
        embedded = self.embedding(x)            # (batch, seq_len, embed_dim)
        output, (hidden, cell) = self.rnn(embedded)
        # Concatenar últimos estados forward y backward
        hidden_final = torch.cat((hidden[-2,:,:], hidden[-1,:,:]), dim=1)
        out = self.fc(self.dropout(hidden_final))
        return out
    \end{lstlisting}
\end{frame}

\begin{frame}[fragile]{Manejo de padding en RNN}
    \begin{lstlisting}
from torch.nn.utils.rnn import pack_padded_sequence, pad_packed_sequence

def forward(self, x, lengths):
    embedded = self.embedding(x)
    packed = pack_padded_sequence(embedded, lengths.cpu(),
                                   batch_first=True, enforce_sorted=False)
    packed_output, (hidden, cell) = self.rnn(packed)
    output, _ = pad_packed_sequence(packed_output, batch_first=True)
    # Usar hidden o output con máscara...
    \end{lstlisting}
    % Basado en Pregunta 7
\end{frame}

% ============================================================================
% SECCIÓN 8: APLICACIONES PRÁCTICAS
% ============================================================================
\section{Aplicaciones Prácticas}

\begin{frame}{Clasificación de reseñas (análisis de sentimiento)}
    \begin{itemize}
        \item Entrada: secuencia de palabras.
        \item LSTM bidireccional + promedio de estados o último estado.
        \item Clasificación binaria (positivo/negativo) o multiclase.
    \end{itemize}
    % Basado en Preguntas 1 y 10
\end{frame}

\begin{frame}{Reconocimiento de Entidades Nombradas (NER)}
    \begin{itemize}
        \item Etiquetado por token (B-PER, I-PER, O, etc.).
        \item LSTM bidireccional + capa lineal + softmax en cada paso.
        \item Pérdida: cross-entropy por token.
    \end{itemize}
    % Basado en Preguntas 6 y 13
\end{frame}

\begin{frame}{Modelado de lenguaje / generación de texto}
    \begin{itemize}
        \item Predecir la siguiente palabra dado el contexto.
        \item LSTM + softmax sobre vocabulario.
        \item Evaluación: perplejidad (perplexity).
        \item Decodificación: greedy, sampling con temperatura.
    \end{itemize}
    % Basado en Preguntas 8 y 9
\end{frame}

\begin{frame}{Traducción automática (seq2seq)}
    \begin{itemize}
        \item Encoder LSTM: codifica frase origen en un vector de contexto.
        \item Decoder LSTM: genera frase destino a partir del contexto.
        \item Limitación: vector de contexto fijo para frases largas.
        \item Mejora: atención (próxima sesión).
    \end{itemize}
    % Basado en Pregunta 16
\end{frame}

% ============================================================================
% SECCIÓN 9: CASO INTEGRADO - ANÁLISIS DE RESEÑAS HOTELERAS
% ============================================================================
\section{Caso Integrado: Análisis de Reseñas Hoteleras}

\begin{frame}{Problema}
    \begin{itemize}
        \item Cadena hotelera recibe miles de reseñas.
        \item Objetivos:
        \begin{enumerate}
            \item Clasificar sentimiento global (positivo/negativo/neutral).
            \item Extraer aspectos (limpieza, ubicación, personal, precio) y su sentimiento asociado.
        \end{enumerate}
    \end{itemize}
    % Basado en Pregunta 20
\end{frame}

\begin{frame}{Solución propuesta}
    \begin{columns}[t]
        \column{0.5\textwidth}
        \textbf{Tarea 1: Clasificación global}
        \begin{itemize}
            \item Embedding + LSTM bidireccional.
            \item Promedio de estados ocultos.
            \item Capa fully-connected + softmax.
        \end{itemize}
        \column{0.5\textwidth}
        \textbf{Tarea 2: Extracción de aspectos}
        \begin{itemize}
            \item Misma capa de embedding y LSTM.
            \item Salida por token (etiquetas BIO).
            \item Clasificador de sentimiento por aspecto.
        \end{itemize}
    \end{columns}
\end{frame}

\begin{frame}{Multi-task learning}
    \begin{itemize}
        \item Arquitectura compartida (embedding + LSTM) con dos cabezas.
        \item Ventaja: representaciones compartidas mejoran ambas tareas.
        \item Desventaja: mayor complejidad de entrenamiento.
        \item Métricas: F1 para clasificación, F1 a nivel de entidad para aspectos.
    \end{itemize}
\end{frame}

% ============================================================================
% SECCIÓN 10: PREPARACIÓN PARA ENTREVISTAS
% ============================================================================
\section{Preparación para Entrevistas}

\begin{frame}{Preguntas típicas}
    \begin{itemize}
        \item ¿Por qué una red feed-forward con average pooling falla en distinguir oraciones con el mismo vocabulario pero diferente orden?
        \item Explica el problema del gradiente desvaneciente en RNN y cómo lo solucionan LSTM/GRU.
        \item Escribe las ecuaciones de una celda LSTM y explica cada puerta.
        \item Compara LSTM y GRU. ¿Cuándo usarías cada una?
        \item ¿Qué es una RNN bidireccional y para qué sirve?
        \item ¿Cómo manejarías secuencias de longitudes variables en un lote?
    \end{itemize}
    % Basado en solucionario
\end{frame}

\begin{frame}{Preguntas avanzadas}
    \begin{itemize}
        \item ¿Cómo funciona backpropagation through time (BPTT)?
        \item ¿Qué es el truncamiento en BPTT y por qué es necesario?
        \item ¿Qué técnicas de regularización usarías en una LSTM?
        \item ¿Cómo inicializarías los embeddings y los pesos de las puertas?
        \item ¿Qué es perplejidad y cómo se interpreta?
    \end{itemize}
    % Basado en Preguntas 5, 9, 18, 19
\end{frame}

% ============================================================================
% SECCIÓN 11: CONEXIÓN CON EL RESTO DEL CURSO
% ============================================================================
\section{Conexión con el resto del curso}

\begin{frame}{Próximos pasos}
    \begin{itemize}
        \item \textbf{Sesión 4 (Embeddings)}: base para capas de embedding.
        \item \textbf{Sesión 5 (Modelos clásicos)}: superados por RNN en secuencias.
        \item \textbf{Sesión 7 (Atención y Transformers)}: superan limitaciones de RNN (paralelización, dependencias largas).
    \end{itemize}
\end{frame}

% ============================================================================
% SECCIÓN 12: RESUMEN
% ============================================================================
\section{Resumen}

\begin{frame}{Lo que hemos aprendido}
    \begin{exampleblock}{Conceptos clave}
        \begin{itemize}
            \item \textbf{RNN vanilla}: procesan secuencias, sufren gradiente desvaneciente.
            \item \textbf{LSTM}: puertas de olvido, entrada, salida; estado de celda preserva memoria.
            \item \textbf{GRU}: simplificación de LSTM, menos parámetros, más rápido.
            \item \textbf{BPTT}: algoritmo de entrenamiento para RNN.
            \item \textbf{RNN bidireccionales}: contexto completo, ideales para NER.
            \item \textbf{Regularización}: dropout, weight decay, clipping.
            \item \textbf{Aplicaciones}: clasificación, NER, modelado de lenguaje, traducción.
            \item \textbf{Caso integrado}: análisis de reseñas con multi-task.
        \end{itemize}
    \end{exampleblock}
\end{frame}

% --- Diapositiva final ---
\begin{frame}
    \centering
    \Huge \textbf{¡Gracias!}

\end{frame}

\end{document}